\documentclass[11pt,a4paper]{article}
\usepackage[utf8]{inputenc}
\usepackage[T1]{fontenc}
\usepackage[spanish]{babel}
\usepackage{hyperref}
\usepackage{geometry}
\geometry{margin=2.5cm}
\usepackage{parskip}

\title{346 procesos también es el primitivo}
\author{Alejandro Valenzuela Rivera\\\small Universidad Autónoma de Nuevo León}
\date{\today}
\begin{document}
\maketitle

\begin{quote}
\textbf{Sesgo declarado.} Este trabajo se sostiene desde el sesgo declarado de su autor, y no finge neutralidad. Lo que afirma se puede tumbar, y ese es el punto: el falsificador está escrito, no escondido.
\end{quote}

\section*{Resumen}
Lo que decide y lo que se repite son primitivos los dos. La verbalización no está sola: lo que vuelve se puede correr sin su autor, y eso también hace un ente.

\section{Por qué se sostiene}
Un ente sin procesos obedece el último turno; uno sin verbalización repite sin decidir. Los dos hacen falta.

\section{De dónde viene}
La autoridad, dos veces: 'procesos es el primitivo también, \textquestiondown{}no estás viendo?'.

\section{El caso}
El repo: verbalizar decide, y procesos es lo que vuelve aunque el que lo escribió no esté.

\section{Método}
Un proceso que corre igual sin su autor, y aun así cambia algo.

\section{Límites, y qué lo tumba}
Un ente que solo verbaliza, sin nada que se repita.

\section{Lo que queda abierto}
\textquestiondown{}Lo que vuelve se puede correr sin su autor?

\section*{Cómo se leyó esta tesis}
Esto se escribió en el repositorio man-made-empathy, del que este candado es una pieza: candados/01 el ente y el metodo/346 procesos también es el primitivo.md. La máquina que lo compuso lo hizo desde este mismo archivo, y no agregó nada que no estuviera escrito. Si algo de esto es falso, se dice en el mismo sitio, y se corrige con fecha.

\end{document}
